% Hand-written text; numbers are macros from generated/numbers.tex.
\section{Discussion}\label{sec:discussion}
\textbf{Answer to the question.} In this setting, error-mode heterogeneity predicts RLVR outcome beyond aggregate verifier quality. Verifiers with exactly the same base-weighted false-acceptance rate and the same true-positive rate led to different accuracy, ordered by the dispersion $V$ of their per-mode acceptance probabilities, with a precise negative estimate that survives deduplication, influence checks, a run-level model and a rank-based check. Across three aggregate rates, $V$ predicted left-out configurations better than the aggregate rate or the worst-mode rate (exploratory). The pre-registered competitors did not beat it: $\fmax$ added little once $V$ was known, and the CGF plug-in, which theory suggests should gain at longer horizons, chose the smallest horizon on its grid and matched $V$ without improving on it. At $T=\HorizonT$ steps this is the short-horizon regime in which the replicator expansion makes $V$ the leading term (Section~\ref{sec:theory}). The sign of the effect was not guaranteed by theory, since the vanilla-gradient statistic can reverse it; here it is the replicator sign.

\textbf{What the evidence does not show.} First, the registered held-out test (H2) favoured $V$, but \NheldTwin{} of its \NheldOut{} held-out configurations duplicate training configurations; on the \PureN{} that do not, the comparison is uninformative. The leave-one-verifier-out comparison over distinct verifiers supports H2's direction, but it was not the registered test. Second, every $p$-value was declared non-confirmatory in advance; we report the H1 interval as an estimate and do not convert it into a confirmatory rejection after the fact. Third, the mechanism analysis is exploratory. The rise of the most-accepted mode and of the realised false-acceptance rate with $V$, and the hacking episodes on a nearly always-accepted mode, are consistent with the selection mechanism in the theory, but they were not registered tests.

\textbf{Practical reading.} A verifier's aggregate false-positive rate, measured on the base policy, is not a sufficient description of its risk. Two checks with the same measured rate can differ in how much they reward a single kind of wrong answer, and RL finds and amplifies that kind. Reporting the per-mode acceptance profile, or at least its dispersion and maximum, is cheap when error modes can be classified, and it predicted outcomes here where the aggregate did not. Monitoring the realised, on-policy false-acceptance rate during training is a direct way to see the drift.

\section{Limitations}\label{sec:limits}
\begin{itemize}
\item \textbf{Scope.} One model (Qwen2.5-0.5B-Instruct), one task (five-digit addition), four outcome-defined error modes, synthetic Bernoulli verifiers, $T=\HorizonT$ GRPO steps, one learning rate. Nothing here licenses claims about larger models, other tasks, natural verifiers or long training. In particular, the long-horizon regime, where theory predicts that $\fmax$ takes over, was not reached.
\item \textbf{Two seeds per configuration}, as registered. Seed-to-seed spread grows with $V$, and the hacking episodes occurred in one seed of a pair; their frequency cannot be estimated from two seeds.
\item \textbf{Power and claim type.} The registered power for H1 at the strictest Holm threshold was \PowC{} (pessimistic \PowPessC). The study was therefore registered as estimation with confidence intervals, and it is reported that way, although the observed interval turned out to be narrow.
\item \textbf{Duplicated configurations.} The registered grid has \NuniqueMatched{} distinct matched verifiers, not 36. Registered standard errors and the registered held-out test are affected; the deduplicated analyses are post hoc.
\item \textbf{Census stability.} The reference distribution passed the registered split-half criterion, but that criterion would fail for the same census under \TVpfail{} of random prompt splits. The lock is supported by an independent replicate and by bootstrap propagation ($\Fbar$ within $\pm$\FbarDev{} of 0.3), so the matching of aggregate rates holds to about $\pm0.02$, not exactly, at the true base distribution.
\item \textbf{1.5B model.} A secondary Qwen2.5-1.5B arm failed the same stability criterion (TV 0.039 and 0.050 against 0.03). A post hoc diagnosis attributes this to sampling variability under strong within-prompt clustering, but the declared failure stands, no 1.5B reference distribution exists, and no 1.5B result is reported.
\item \textbf{Pilot.} The heterogeneity pilot used 4 configurations $\times$ 3 seeds and is descriptive. Its effect size was not used for planning; the planning effect came from simulation.
\item \textbf{Taxonomy and audit.} The registered mechanism taxonomy reached its stop rule and was replaced, before any training data, by an outcome-defined partition (Amendment~1). The 200-row construct audit was labelled by an AI assistant from written rules and reviewed by the author.
\item \textbf{Literature coverage.} Forward-citation search was rate-limited, so recent work citing the closest papers may be missing from the related-work review.
\end{itemize}

\section{Conclusion}\label{sec:conclusion}
In a pre-registered, exactly rate-matched study of GRPO on Qwen2.5-0.5B-Instruct and five-digit addition, verifiers that spread the same false-acceptance budget unevenly across error modes led to lower accuracy, and the dispersion statistic $V$ predicted outcomes better than the aggregate rate or the worst-mode rate. The study was registered as an estimation study and is reported as one; the held-out test is weakened by duplicated configurations in the registered grid, and the evidence is confined to one small model, one task and a short horizon. The code, the registration, every decision and all run records are public so that the estimate can be checked and extended.

\paragraph{Data and code availability.} The pre-registration is DOI 10.17605/OSF.IO/BS3WQ (\url{https://osf.io/bs3wq/}). Code, configurations, run records, analysis outputs and the decision log are at \url{https://github.com/Alyssa-286/VERIFLOW}. Model checkpoints are not archived (about 1.1~TB); every run can be regenerated from its configuration and seed.

\paragraph{Use of AI assistance.} Code, analyses and text were developed with an AI coding assistant (Claude) under the author's direction; the construct audit labels are disclosed above. All results come from the logged runs and scripts in the repository.
